%%% MA_13_Growth.tex -- reorganised; frames moved from the old MA_xx files (see comments)

\section{Economic Growth, Innovation, and Resources}

%%% [moved from MA_9_Growth.tex, frame 0]
\begin{frame}{Economic Growth}
    \mytext{
    Economic growth is the growth of value in a microeeconomic sense.

    This means \enquote{more is more}, but it can be that \textit{less of something} is \textit{more value}. Then having \textit{less} of that something corresponds to growth.

    Thus \textit{de-growth} is paradoxical.
    }
\end{frame}

%%% [moved from MA_9_Growth.tex, frame 1]
\begin{frame}{Economic Growth}
    \mytext{
    An economy can grow by several factors, such as:
    \begin{itemize}
        \item growth of population
        \item growth of capital
        \item growth of efficiency:
        \begin{itemize}
            \item more appropriate allocation
            \item increase in knowledge (labour productivity)
            \item technological progress (capital productivity)
            \item land productivity (fertility, drainage etc.)
        \end{itemize}
    \end{itemize}
    }
\end{frame}

%%% [moved from MA_9_Growth.tex, frame 2]
\begin{frame}{Solow-Swan Model}
    \mytext{
        Output $Y$ depends on capital $K$ and labour $L$. 

        Capital is reimbursed by savings and labour is reimbursed by wages.

        Both capital and labour show diminishing returns, thus  $0<\alpha<1$ in a Cobb-Douglas-Production function:

        \begin{equation}
            Y(t) = K(t)^\alpha \cdot L(t)^{1-\alpha}
        \end{equation}
    }
\end{frame}

%%% [moved from MA_9_Growth.tex, frame 3]
\begin{frame}{Solow-Swan Model}
    \mycomment{
        All units are per worker.
        
        $y$ is output
    }
    \showmytikz[1]{../shared/tikz/tikz-solow}
\end{frame}

%%% [moved from MA_9_Growth.tex, frame 4]
\begin{frame}{Solow-Swan Model}
    \mycomment{
        $\delta k$ is depreciation. (The slope is the rate of depreciation; it is not per se $45°$). This is the break-even condition for investment.
    }
    \showmytikz[1,2]{../shared/tikz/tikz-solow}
\end{frame}

%%% [moved from MA_9_Growth.tex, frame 5]
\begin{frame}{Solow-Swan Model}
    \mycomment{
        $sy_t$ is investment. This equals the part of the income which is saved. The rate is diminishing because of the diminishing returns of $y(k)$. 
    }
    \showmytikz[1,2,3]{../shared/tikz/tikz-solow}
\end{frame}

%%% [moved from MA_9_Growth.tex, frame 6]
\begin{frame}{Solow-Swan Model}
    \mycomment{
        With this in place, we see a steady state income of $y^*$. 
        
        At the steady state, investment will just replace depreciation ($E$). Then growth only happens because of technological progress (the famous \enquote{Solow Residual}) and/or population growth.
    }
    \showmytikz[1,2,3,4,5]{../shared/tikz/tikz-solow}
\end{frame}


%%% TODO (new content): convergence; growth accounting; institutions

%%% [moved from MA_9_Growth.tex, frame 16]
\begin{frame}{Endogenous Growth Theory: Education and Research}
    \mytext{
    State should finance education and fundamental research
    \begin{itemize}
        \item knowledge is non-rival
        \item knowledge keeps compounding (?)
        \item thus not diminishing returns
        \item[$\rightarrow$] constant or increasing returns
        \item[$\rightarrow$] exponential growth        
    \end{itemize}
    }
    \mycomment{
    The argument for education and knowledge as ultimate driver of economic growth was formalised by \noun{Lucas} and \noun{Romer}. There is critique, too: 
    \begin{enumerate}
        \item explains inexplicable variables by further inexplicable variables (\noun{Krugman})
        \item why should knowledge not show diminishing returns?
        \item there should be / would be some sort of plateau anyway (e.g. physical boundaries of computer capacity on the atomic level)
    \end{enumerate}
    }
\end{frame}

\begin{frame}{Distribution of Growth}
    \mycomment{
    Recently, the term \enquote{K-shaped recovery} (or K-shaped economy) has become popular. The idea is that growth (recovery after Covid) has been unevenly distributed, with the top (and bottom) recovering faster.
    }
    \showmytikz[1]{../shared/tikz/tikz-k-shape}
\end{frame}

\begin{frame}{Distribution of Growth}
    \mycomment{
    There has also been talk of n \enquote{E-shaped} recovery.
    }
    \showmytikz[2]{../shared/tikz/tikz-k-shape}
\end{frame}

%%% [moved from MA_9_Growth.tex, frame 17]
\begin{frame}{Innovation Economics}
    \mytext{
    \showmytikz{../shared/tikz/tikz-schumpeter}
    \begin{itemize}
        \item Innovation can lead to growth
        \item innovation devalues older technology / investment
    \end{itemize}
    }
\end{frame}

%%% [moved from MA_10_Cycles.tex, frame 9]
\begin{frame}{Creative Descruction}
    \showmytikz{../shared/tikz/tikz-schumpeter}
    \mycomment{
    The new technology, shown in green, overtakes the old technology, shown in red. Total output increases seamlessly since the new technology is more productive than the old one.
    }
    
\end{frame}

%%% [moved from MA_9_Growth.tex, frame 14]
\begin{frame}{Triple Helix Model}
    \showmytikz[1]{../shared/tikz/tikz-triplehelix}
\end{frame}

%%% [moved from MA_9_Growth.tex, frame 15]
\begin{frame}{Intellectual Property}
    \mytext{
    Specific technical solutions to specific technological problems can be patented.
    \begin{itemize}
        \item spending on research only pays out when findings are protected.
        \item[$\rightarrow$] argument for strong and long IP protection
        \item research only pays out when opportunity exists to use findings
        \item[$\rightarrow$] argument for weaker IP protection
        \item invention-to-market length / funding
        \item technological readiness concept
        \item market readiness concept
    \end{itemize}
    }
\end{frame}

%%% [moved from MA_9_Growth.tex, frame 7]
\begin{frame}{Club of Rome / Meadow: World3}
    \showmytikz[1]{../shared/tikz/tikz-world3}
\end{frame}

%%% [moved from MA_9_Growth.tex, frame 8]
\begin{frame}{Club of Rome / Meadow: World3}
    \showmytikz[1,2]{../shared/tikz/tikz-world3}
\end{frame}

%%% [moved from MA_9_Growth.tex, frame 9]
\begin{frame}{Club of Rome / Meadow: World3}
    \showmytikz[1,2,3]{../shared/tikz/tikz-world3}
\end{frame}

%%% [moved from MA_9_Growth.tex, frame 10]
\begin{frame}{Club of Rome / Meadow: World3}
    \showmytikz[1,2,3,4]{../shared/tikz/tikz-world3}
\end{frame}

%%% [moved from MA_9_Growth.tex, frame 11]
\begin{frame}{Club of Rome / Meadow: World3}
    \showmytikz[1,2,3,4,5]{../shared/tikz/tikz-world3}
\end{frame}

%%% [moved from MA_9_Growth.tex, frame 12]
\begin{frame}{Club of Rome / Meadow: World3}
    \showmytikz[1,2,3,4,5,6]{../shared/tikz/tikz-world3}
\end{frame}

%%% [moved from MA_9_Growth.tex, frame 13]
\begin{frame}{Resources Economics}
    \showmytikz[1,2]{../shared/tikz/tikz-hotellingpath}
    \mycomment{
    As resources become more rare, their price increases on the Hotelling path.

    This makes more expensive extraction viable, e.g. offshore drilling.
    
    It also makes substitutes and efficiency gain more attractive.

    There may be \enquote{peak oil,} but there is no \enquote{last drop of oil}.

    Will there be cheaper synthetic alternatives? It could be that the price of alternatives actually depends on / follows  the price of the original product (per Merit Order effect).
    }
\end{frame}

